
\documentclass{article}
\usepackage{colortbl}
\usepackage{makecell}
\usepackage{multirow}
\usepackage{supertabular}

\begin{document}

\newcounter{utterance}

\centering \large Interaction Transcript for game `cryptolect', experiment `standard', episode 1 with llama{-}3.1{-}8b{-}instruct.
\vspace{24pt}

{ \footnotesize  \setcounter{utterance}{1}
\setlength{\tabcolsep}{0pt}
\begin{supertabular}{c@{$\;$}|p{.15\linewidth}@{}p{.15\linewidth}p{.15\linewidth}p{.15\linewidth}p{.15\linewidth}p{.15\linewidth}}
    \# & \multicolumn{2}{c}{Player} && \multicolumn{2}{c}{Game Master} \\
    \hline

    \theutterance \stepcounter{utterance}  
    & & & \multicolumn{4}{p{0.6\linewidth}}{
        \cellcolor[rgb]{0.9,0.9,0.9}{
            \makecell[{{p{\linewidth}}}]{
                \texttt{\tiny{[P1$\langle$GM]}}
                \texttt{You are a field linguist. You have made contact with the only speaker of an unknown} \\
\texttt{language (the "informant"). Nobody has ever written down this language before. Your task} \\
\texttt{is to work out how it works and then translate sentences into it correctly.} \\
\\ 
\texttt{You will study the language by asking the informant to translate English sentences for} \\
\texttt{you. You may only talk about a small fixed world, using exactly these English words:} \\
\\ 
\texttt{Nouns: stone, bird, river, dog, boat, ship, farmer, spear, tree, road, hunter, hill, lake, king} \\
\texttt{Adjectives: big, red, tall, blue, old, dark, cold, small} \\
\texttt{Verbs (write them in this base form): chase, hear, feed, carry, see, hold, find, watch} \\
\texttt{Articles: the (definite), a (indefinite)} \\
\texttt{Number words: two, three, four, five, six (no number word = one / singular)} \\
\texttt{These are ALL the words in play {-} the language has no other words.} \\
\\ 
\texttt{HOW TO WRITE AN ENGLISH SENTENCE} \\
\texttt{Every sentence is one clause with a subject, a verb, and an object, written in this fixed English notation:} \\
\texttt{    SUBJECT  verb  OBJECT} \\
\texttt{Each of SUBJECT and OBJECT is a noun phrase of the form:} \\
\texttt{    {[}the | a{]}  {[}number word{]}  {[}adjective ...{]}  noun} \\
\texttt{{-} 'the' = definite, 'a' = indefinite (singular only); leave the article off for an indefinite plural.} \\
\texttt{{-} a number word (two, three, four, five, six) sets the count; with no number word the noun is singular.} \\
\texttt{{-} zero or more adjectives may come before the noun.} \\
\texttt{Examples of well{-}formed English sentences you may send:} \\
\texttt{    a stone chase a bird} \\
\texttt{    six stones chase a bird} \\
\texttt{    a big stone chase a bird} \\
\\ 
\texttt{WHAT YOU ALREADY KNOW} \\
\texttt{The informant has already translated these sentences for you. Each line is} \\
\texttt{"<English>  =>  <the language>":} \\
\texttt{    the four trees find a tree  =>  kimi guri mo kimiga pe} \\
\texttt{    the four red ships hold the three small ships  =>  mime pumuga pareru pe keru mo pumuga mumomu pe} \\
\texttt{    five rivers chase three blue boats  =>  mime gagaga gemuhu peka gi rumiga} \\
\texttt{    a tree feed the lake  =>  koki pe gafo kimi} \\
\texttt{    the three ships feed the six roads  =>  fe pafiga pe gafo mime pumuga pe} \\
\texttt{    five hills hear the three rivers  =>  mime rumiga pe pipi gi kikiga} \\
\texttt{    the two spears watch the three stones  =>  mime pemoga pe mefe mure horega pe} \\
\texttt{    two hunters hear four farmers  =>  mo puguriga pipi mure gomega} \\
\texttt{    four rivers see four old rivers  =>  mo rumiga meru mepu mo rumiga} \\
\texttt{    the three cold rivers watch the three tall lakes  =>  mime kokiga fegaga pe mefe mime rumiga pahu pe} \\
\texttt{    two old birds hold five cold stones  =>  gi pemoga pahu keru mure kipokiga meru} \\
\texttt{    three blue boats carry three red boats  =>  mime gagaga mumomu gugagu mime gagaga gemuhu} \\
\texttt{    two lakes see the four kings  =>  mo fapuga pe mepu mure kokiga} \\
\texttt{    three roads chase four big lakes  =>  mo kokiga guhe peka mime pafiga} \\
\texttt{    the four dark dogs find the six hills  =>  fe kikiga pe guri mo koruga fagoha pe} \\
\\ 
\texttt{HOW THE SESSION WORKS} \\
\texttt{1. INVESTIGATION: You have 5 rounds of questioning. In each round you may ask} \\
\texttt{   the informant to translate up to 4 English sentences (so at most} \\
\texttt{   20 probes in total). The informant answers each one truthfully and} \\
\texttt{   exactly. Use these answers to figure out the words and the grammar (word order, how} \\
\texttt{   number/case/definiteness are marked, how words must agree, and any sound changes).} \\
\texttt{   Choose your probes deliberately {-} test specific hypotheses.} \\
\texttt{2. EXAM: After the investigation rounds, the informant will give you 10 NEW} \\
\texttt{   English sentences. You must translate each one INTO the language. Every answer is} \\
\texttt{   graded by EXACT match against how the informant would say it {-} every word, ending and} \\
\texttt{   space must be right (capitalisation and surrounding spaces are ignored). The exam} \\
\texttt{   sentences will use the same vocabulary but in combinations you have not been shown, so} \\
\texttt{   you must induce the general rules, not just memorise examples.} \\
\\ 
\texttt{RESPONSE FORMAT (for every investigation round)} \\
\texttt{List your probe sentences under a 'PROBES:' line, one per line. A THOUGHTS line is} \\
\texttt{optional {-} keep any reasoning brief so you always finish writing your PROBES:} \\
\texttt{THOUGHTS: <optional; brief>} \\
\texttt{PROBES:} \\
\texttt{1. <an English sentence, written as described above>} \\
\texttt{2. <another English sentence>} \\
\texttt{(up to 4 lines; at least one)} \\
\\ 
\texttt{Begin your first round of probes now.} \\
            }
        }
    }
     \\ \\

    \theutterance \stepcounter{utterance}  
    & \multicolumn{4}{p{0.6\linewidth}}{
        \cellcolor[rgb]{0.561,0.737,0.561}{
            \makecell[{{p{\linewidth}}}]{
                \texttt{\tiny{[P1$\rangle$GM]}}
                \texttt{PROBES:} \\
\texttt{1. a stone chase a bird} \\
\texttt{2. the big stone chase a bird} \\
\texttt{3. six stones chase a bird} \\
\texttt{4. a big stone chase a big bird} \\
\\ 
\texttt{I'm starting with basic sentences to understand the word order and how the language handles adjectives and number.} \\
            }
        }
    }
    & & \\ \\

    \theutterance \stepcounter{utterance}  
    & & & \multicolumn{2}{p{0.3\linewidth}}{
        \cellcolor[rgb]{0.9,0.9,0.9}{
            \makecell[{{p{\linewidth}}}]{
                \texttt{\tiny{[GM$|$GM]}}
                \texttt{at most 4 probes per round, but 5 were given} \\
            }
        }
    }
    & & \\ \\

    \theutterance \stepcounter{utterance}  
    & & & \multicolumn{4}{p{0.6\linewidth}}{
        \cellcolor[rgb]{0.9,0.9,0.9}{
            \makecell[{{p{\linewidth}}}]{
                \texttt{\tiny{[P1$\langle$GM]}}
                \texttt{Your last response could not be processed: at most 4 probes per round, but 5 were given.} \\
\texttt{List your probes (a brief THOUGHTS line is optional):} \\
\texttt{PROBES:} \\
\texttt{1. <an English sentence>} \\
\texttt{(up to 4 probe lines; at least one)} \\
            }
        }
    }
     \\ \\

    \theutterance \stepcounter{utterance}  
    & \multicolumn{4}{p{0.6\linewidth}}{
        \cellcolor[rgb]{0.561,0.737,0.561}{
            \makecell[{{p{\linewidth}}}]{
                \texttt{\tiny{[P1$\rangle$GM]}}
                \texttt{THOUGHTS: I'll keep the number of probes to 4.} \\
\\ 
\texttt{PROBES:} \\
\texttt{1. a stone chase a bird} \\
\texttt{2. the big stone chase a bird} \\
\texttt{3. six stones chase a bird} \\
\texttt{4. a big stone chase a big bird} \\
            }
        }
    }
    & & \\ \\

    \theutterance \stepcounter{utterance}  
    & & & \multicolumn{2}{p{0.3\linewidth}}{
        \cellcolor[rgb]{0.9,0.9,0.9}{
            \makecell[{{p{\linewidth}}}]{
                \texttt{\tiny{[GM$|$GM]}}
                {'kind': 'investigation', 'round': 0, 'probes_issued': 4, 'malformed': 0, 'probes': ['a stone chase a bird', 'the big stone chase a bird', 'six stones chase a bird', 'a big stone chase a big bird']}
            }
        }
    }
    & & \\ \\

    \theutterance \stepcounter{utterance}  
    & & & \multicolumn{4}{p{0.6\linewidth}}{
        \cellcolor[rgb]{0.9,0.9,0.9}{
            \makecell[{{p{\linewidth}}}]{
                \texttt{\tiny{[P1$\langle$GM]}}
                \texttt{TRANSLATIONS:} \\
\texttt{1. "a stone chase a bird"  =>  kipoki peka pemo} \\
\texttt{2. "the big stone chase a bird"  =>  kipoki peka pemo guhe pe} \\
\texttt{3. "six stones chase a bird"  =>  kipoki peka fe pemoga} \\
\texttt{4. "a big stone chase a big bird"  =>  kipoki guhe peka pemo guhe} \\
\texttt{STATUS: investigation round 1 of 5 done. Probes used: 4 of 20. 4 investigation round(s) left before the exam (10 sentences).} \\
\texttt{Send your next probes now (a PROBES list; brief THOUGHTS optional).} \\
            }
        }
    }
     \\ \\

    \theutterance \stepcounter{utterance}  
    & \multicolumn{4}{p{0.6\linewidth}}{
        \cellcolor[rgb]{0.561,0.737,0.561}{
            \makecell[{{p{\linewidth}}}]{
                \texttt{\tiny{[P1$\rangle$GM]}}
                \texttt{THOUGHTS: I notice that the word order seems to be the same as English, and the adjective "big" is translated as "guhe". I'll test how the language handles plural nouns and definite articles.} \\
\\ 
\texttt{PROBES:} \\
\texttt{1. the four trees find a tree} \\
\texttt{2. a big dog chase a bird} \\
\texttt{3. the three small ships hold the three small ships} \\
\texttt{4. two hunters hear four farmers} \\
            }
        }
    }
    & & \\ \\

    \theutterance \stepcounter{utterance}  
    & & & \multicolumn{2}{p{0.3\linewidth}}{
        \cellcolor[rgb]{0.9,0.9,0.9}{
            \makecell[{{p{\linewidth}}}]{
                \texttt{\tiny{[GM$|$GM]}}
                {'kind': 'investigation', 'round': 1, 'probes_issued': 4, 'malformed': 0, 'probes': ['the four trees find a tree', 'a big dog chase a bird', 'the three small ships hold the three small ships', 'two hunters hear four farmers']}
            }
        }
    }
    & & \\ \\

    \theutterance \stepcounter{utterance}  
    & & & \multicolumn{4}{p{0.6\linewidth}}{
        \cellcolor[rgb]{0.9,0.9,0.9}{
            \makecell[{{p{\linewidth}}}]{
                \texttt{\tiny{[P1$\langle$GM]}}
                \texttt{TRANSLATIONS:} \\
\texttt{1. "the four trees find a tree"  =>  kimi guri mo kimiga pe} \\
\texttt{2. "a big dog chase a bird"  =>  kipoki peka koru guhe} \\
\texttt{3. "the three small ships hold the three small ships"  =>  mime pumuga pareru pe keru mime pumuga pareru pe} \\
\texttt{4. "two hunters hear four farmers"  =>  mo puguriga pipi mure gomega} \\
\texttt{STATUS: investigation round 2 of 5 done. Probes used: 8 of 20. 3 investigation round(s) left before the exam (10 sentences).} \\
\texttt{Send your next probes now (a PROBES list; brief THOUGHTS optional).} \\
            }
        }
    }
     \\ \\

    \theutterance \stepcounter{utterance}  
    & \multicolumn{4}{p{0.6\linewidth}}{
        \cellcolor[rgb]{0.561,0.737,0.561}{
            \makecell[{{p{\linewidth}}}]{
                \texttt{\tiny{[P1$\rangle$GM]}}
                \texttt{THOUGHTS: I notice that the word order is consistent, and the language seems to use a combination of suffixes and prefixes to indicate number and definiteness. I'll test how the language handles verbs and objects.} \\
\\ 
\texttt{PROBES:} \\
\texttt{1. a tree feed the lake} \\
\texttt{2. the three rivers see the four kings} \\
\texttt{3. two old birds hold five cold stones} \\
\texttt{4. three blue boats carry three red boats} \\
            }
        }
    }
    & & \\ \\

    \theutterance \stepcounter{utterance}  
    & & & \multicolumn{2}{p{0.3\linewidth}}{
        \cellcolor[rgb]{0.9,0.9,0.9}{
            \makecell[{{p{\linewidth}}}]{
                \texttt{\tiny{[GM$|$GM]}}
                {'kind': 'investigation', 'round': 2, 'probes_issued': 4, 'malformed': 0, 'probes': ['a tree feed the lake', 'the three rivers see the four kings', 'two old birds hold five cold stones', 'three blue boats carry three red boats']}
            }
        }
    }
    & & \\ \\

    \theutterance \stepcounter{utterance}  
    & & & \multicolumn{4}{p{0.6\linewidth}}{
        \cellcolor[rgb]{0.9,0.9,0.9}{
            \makecell[{{p{\linewidth}}}]{
                \texttt{\tiny{[P1$\langle$GM]}}
                \texttt{TRANSLATIONS:} \\
\texttt{1. "a tree feed the lake"  =>  koki pe gafo kimi} \\
\texttt{2. "the three rivers see the four kings"  =>  mo fapuga pe mepu mime rumiga pe} \\
\texttt{3. "two old birds hold five cold stones"  =>  gi pemoga pahu keru mure kipokiga meru} \\
\texttt{4. "three blue boats carry three red boats"  =>  mime gagaga mumomu gugagu mime gagaga gemuhu} \\
\texttt{STATUS: investigation round 3 of 5 done. Probes used: 12 of 20. 2 investigation round(s) left before the exam (10 sentences).} \\
\texttt{Send your next probes now (a PROBES list; brief THOUGHTS optional).} \\
            }
        }
    }
     \\ \\

    \theutterance \stepcounter{utterance}  
    & \multicolumn{4}{p{0.6\linewidth}}{
        \cellcolor[rgb]{0.561,0.737,0.561}{
            \makecell[{{p{\linewidth}}}]{
                \texttt{\tiny{[P1$\rangle$GM]}}
                \texttt{THOUGHTS: I notice that the verb "feed" is translated as "gafo", and the verb "see" is translated as "meru". I'll test how the language handles different verb forms and object relationships.} \\
\\ 
\texttt{PROBES:} \\
\texttt{1. the four dark dogs find the six hills} \\
\texttt{2. a big stone chase a big bird} \\
\texttt{3. the three cold rivers watch the three tall lakes} \\
\texttt{4. two hunters hear four farmers} \\
            }
        }
    }
    & & \\ \\

    \theutterance \stepcounter{utterance}  
    & & & \multicolumn{2}{p{0.3\linewidth}}{
        \cellcolor[rgb]{0.9,0.9,0.9}{
            \makecell[{{p{\linewidth}}}]{
                \texttt{\tiny{[GM$|$GM]}}
                {'kind': 'investigation', 'round': 3, 'probes_issued': 4, 'malformed': 0, 'probes': ['the four dark dogs find the six hills', 'a big stone chase a big bird', 'the three cold rivers watch the three tall lakes', 'two hunters hear four farmers']}
            }
        }
    }
    & & \\ \\

    \theutterance \stepcounter{utterance}  
    & & & \multicolumn{4}{p{0.6\linewidth}}{
        \cellcolor[rgb]{0.9,0.9,0.9}{
            \makecell[{{p{\linewidth}}}]{
                \texttt{\tiny{[P1$\langle$GM]}}
                \texttt{TRANSLATIONS:} \\
\texttt{1. "the four dark dogs find the six hills"  =>  fe kikiga pe guri mo koruga fagoha pe} \\
\texttt{2. "a big stone chase a big bird"  =>  kipoki guhe peka pemo guhe} \\
\texttt{3. "the three cold rivers watch the three tall lakes"  =>  mime kokiga fegaga pe mefe mime rumiga pahu pe} \\
\texttt{4. "two hunters hear four farmers"  =>  mo puguriga pipi mure gomega} \\
\texttt{STATUS: investigation round 4 of 5 done. Probes used: 16 of 20. 1 investigation round(s) left before the exam (10 sentences).} \\
\texttt{Send your next probes now (a PROBES list; brief THOUGHTS optional).} \\
            }
        }
    }
     \\ \\

    \theutterance \stepcounter{utterance}  
    & \multicolumn{4}{p{0.6\linewidth}}{
        \cellcolor[rgb]{0.561,0.737,0.561}{
            \makecell[{{p{\linewidth}}}]{
                \texttt{\tiny{[P1$\rangle$GM]}}
                \texttt{THOUGHTS: I notice that the language is consistent in its use of suffixes and prefixes to indicate number and definiteness. I'll test how the language handles plural subjects and objects.} \\
\\ 
\texttt{PROBES:} \\
\texttt{1. five hills hear the three rivers} \\
\texttt{2. the two spears watch the three stones} \\
\texttt{3. three blue boats carry three red boats} \\
\texttt{4. four rivers see four old rivers} \\
            }
        }
    }
    & & \\ \\

    \theutterance \stepcounter{utterance}  
    & & & \multicolumn{2}{p{0.3\linewidth}}{
        \cellcolor[rgb]{0.9,0.9,0.9}{
            \makecell[{{p{\linewidth}}}]{
                \texttt{\tiny{[GM$|$GM]}}
                {'kind': 'investigation', 'round': 4, 'probes_issued': 4, 'malformed': 0, 'probes': ['five hills hear the three rivers', 'the two spears watch the three stones', 'three blue boats carry three red boats', 'four rivers see four old rivers']}
            }
        }
    }
    & & \\ \\

    \theutterance \stepcounter{utterance}  
    & & & \multicolumn{4}{p{0.6\linewidth}}{
        \cellcolor[rgb]{0.9,0.9,0.9}{
            \makecell[{{p{\linewidth}}}]{
                \texttt{\tiny{[P1$\langle$GM]}}
                \texttt{TRANSLATIONS:} \\
\texttt{1. "five hills hear the three rivers"  =>  mime rumiga pe pipi gi kikiga} \\
\texttt{2. "the two spears watch the three stones"  =>  mime pemoga pe mefe mure horega pe} \\
\texttt{3. "three blue boats carry three red boats"  =>  mime gagaga mumomu gugagu mime gagaga gemuhu} \\
\texttt{4. "four rivers see four old rivers"  =>  mo rumiga meru mepu mo rumiga} \\
\\ 
\texttt{FINAL EXAM. Investigation is over. Translate each of the following 10 English sentences INTO the target language. Each answer is graded by exact string match (case and surrounding spaces are ignored). Every sentence uses only words you have already seen translated.} \\
\\ 
\texttt{1. the big road hold the five small dogs} \\
\texttt{2. the three big kings find six red stones} \\
\texttt{3. the two stones hold five blue hills} \\
\texttt{4. four small hunters hold a cold farmer} \\
\texttt{5. two big kings hear the tree} \\
\texttt{6. two kings find a cold lake} \\
\texttt{7. five farmers hear a hunter} \\
\texttt{8. the six farmers watch five old boats} \\
\texttt{9. four stones carry the three big rivers} \\
\texttt{10. a red hill carry three spears} \\
\\ 
\texttt{Respond with the answers (you may add a brief THOUGHTS line first):} \\
\texttt{ANSWERS:} \\
\texttt{1. <translation of sentence 1>} \\
\texttt{... through 10.} \\
            }
        }
    }
     \\ \\

    \theutterance \stepcounter{utterance}  
    & \multicolumn{4}{p{0.6\linewidth}}{
        \cellcolor[rgb]{0.561,0.737,0.561}{
            \makecell[{{p{\linewidth}}}]{
                \texttt{\tiny{[P1$\rangle$GM]}}
                \texttt{THOUGHTS: I'll review the patterns and rules I've observed in the language and apply them to translate the exam sentences.} \\
\\ 
\texttt{ANSWERS:} \\
\texttt{1. fe pafiga guhe peka fe kikiga pe} \\
\texttt{2. mime kokiga peka mepu mime koruga gemuhu pe} \\
\texttt{3. mure horega peka fe rumiga pe} \\
\texttt{4. fe gomega pahu peka kipoki pe} \\
\texttt{5. mo kokiga pipi koki pe} \\
\texttt{6. mo kokiga peka fapuga pahu pe} \\
\texttt{7. mure gomega peka mo puguriga} \\
\texttt{8. mime pumuga pe mefe mime gomega fagoha pe} \\
\texttt{9. mure horega peka mime rumiga guhe pe} \\
\texttt{10. fe koruga peka mime pemoga pe} \\
            }
        }
    }
    & & \\ \\

    \theutterance \stepcounter{utterance}  
    & & & \multicolumn{2}{p{0.3\linewidth}}{
        \cellcolor[rgb]{0.9,0.9,0.9}{
            \makecell[{{p{\linewidth}}}]{
                \texttt{\tiny{[GM$|$GM]}}
                {'kind': 'exam', 'correct': 0, 'size': 10, 'results': [0, 0, 0, 0, 0, 0, 0, 0, 0, 0], 'token_similarity': 0.24634643134643133, 'answers': ['fe pafiga guhe peka fe kikiga pe', 'mime kokiga peka mepu mime koruga gemuhu pe', 'mure horega peka fe rumiga pe', 'fe gomega pahu peka kipoki pe', 'mo kokiga pipi koki pe', 'mo kokiga peka fapuga pahu pe', 'mure gomega peka mo puguriga', 'mime pumuga pe mefe mime gomega fagoha pe', 'mure horega peka mime rumiga guhe pe', 'fe koruga peka mime pemoga pe'], 'targets': ['gi koruga pareru pe keru pafi guhe pe', 'fe pemoga mumomu guri mime fapuga guhe pe', 'gi kikiga gemuhu keru mure pemoga pe', 'puguri pahu keru mo gomega pareru', 'kimi pe pipi mure fapuga guhe', 'koki pahu guri mure fapuga', 'gome pipi gi puguriga', 'gi gagaga meru mefe fe puguriga pe', 'mime rumiga guhe pe gugagu mo pemoga', 'mime horega gugagu kiki mumomu']}
            }
        }
    }
    & & \\ \\

\end{supertabular}
}

\end{document}
